\section{Sistem Kriptografi Merkle-Hellman}

\begin{subcpmk}
  \item Mahasiswa mampu mengimplementasikan prosedur pembangkitan kunci, proses enkripsi, dan proses dekripsi pada sistem Merkle-Hellman Knapsack, serta menganalisis peran pengacakan modular dalam menyembunyikan struktur superincreasing.
\end{subcpmk}

\noindent\textit{Rujukan:} Referensi Knapsack 2026; HAC Bab~8 \cite{hac_online}.

\subsection{Gagasan Utama: Menyembunyikan Struktur}

Kelemahan utama dari barisan superincreasing adalah kemudahannya untuk dipecahkan menggunakan algoritma \textit{greedy}. Oleh karena itu, Ralph Merkle dan Martin Hellman mengusulkan sebuah mekanisme untuk menyembunyikan barisan superincreasing (sebagai kunci privat) menjadi sebuah barisan bilangan bulat acak (sebagai kunci publik) melalui transformasi linear modular.

Kunci publik $P$ tampak seperti barisan acak sehingga penyerang terpaksa menghadapi \textit{Subset Sum Problem} yang NP-complete. Namun, pemilik kunci privat yang mengetahui parameter transformasi dapat mengembalikan $P$ menjadi barisan superincreasing asal, sehingga dekripsi menjadi trivial.

\subsection{Prosedur Operasional Merkle-Hellman}

\subsubsection{Pembangkitan Kunci (Key Generation)}
\begin{enumerate}
    \item \textbf{Memilih Barisan Privat}: Pilih sebuah barisan superincreasing $S = \{s_1, s_2, \dots, s_n\}$.
    \item \textbf{Memilih Parameter Transformasi}: Pilih dua bilangan bulat $m$ dan $q$ sedemikian sehingga:
    \begin{itemize}
        \item $q > \sum_{i=1}^n s_i$ (untuk mencegah overflow saat modulo).
        \item $\gcd(m, q) = 1$ (agar $m$ memiliki invers modulo $q$).
    \end{itemize}
    \item \textbf{Menghitung Kunci Publik}: Hitung setiap elemen kunci publik $p_i$ menggunakan rumus:
    \begin{equation}
      p_i = (s_i \cdot m) \pmod{q}
    \end{equation}
\end{enumerate}
\textbf{Kunci Privat}: $(S, m, q)$ \\
\textbf{Kunci Publik}: $P = \{p_1, p_2, \dots, p_n\}$

\subsubsection{Proses Enkripsi}
Untuk mengenkripsi pesan biner $M = (m_1, m_2, \dots, m_n)$ di mana $m_i \in \{0, 1\}$, pengirim menghitung jumlah terboboti dari kunci publik:
\begin{equation}
  T = \sum_{i=1}^{n} m_i \cdot p_i
\end{equation}
Cipherteks yang dikirimkan adalah nilai skalar $T$.

\subsubsection{Proses Dekripsi}
Penerima menggunakan kunci privat untuk mengembalikan $T$ menjadi bentuk superincreasing $T'$:
\begin{enumerate}
    \item \textbf{Hitung Invers Modular}: Temukan $m^{-1} \pmod{q}$ menggunakan Algoritma Euclidean Diperluas.
    \item \textbf{Transformasi Balik}: Hitung nilai $T'$ sebagai berikut:
    \begin{equation}
      T' = (T \cdot m^{-1}) \pmod{q}
    \end{equation}
    \item \textbf{Ekstraksi Bit}: Karena $T'$ adalah jumlah dari barisan superincreasing $S$, gunakan algoritma \textit{greedy} untuk menentukan bit $m_i$.
\end{enumerate}

\begin{figure}[htbp]
  \centering
  \includegraphics[width=0.8\textwidth]{example-image-a}
  \caption{Alur kerja sistem Merkle-Hellman: Transformasi barisan superincreasing menjadi kunci publik, proses enkripsi melalui penjumlahan subset, dan dekripsi menggunakan invers modular.}
  \label{fig:mh-flow}
\end{figure}

\subsection{Contoh Trace Manual Merkle-Hellman}

Misalkan kita menggunakan parameter berikut:
\begin{itemize}
    \item Barisan privat $S = \{2, 5, 11, 21, 45\}$. (Cek: $2+5+11+21 = 39 < 45 \checkmark$).
    \item Parameter: $m = 7, q = 101$. ($\gcd(7, 101)=1$ dan $101 > 83 \checkmark$).
\end{itemize}

\textbf{1. Pembentukan Kunci Publik $P$}:
\begin{align*}
  p_1 &= (2 \cdot 7) \pmod{101} = 14 \\
  p_2 &= (5 \cdot 7) \pmod{101} = 35 \\
  p_3 &= (11 \cdot 7) \pmod{101} = 77 \\
  p_4 &= (21 \cdot 7) \pmod{101} = 147 \pmod{101} = 46 \\
  p_5 &= (45 \cdot 7) \pmod{101} = 315 \pmod{101} = 12
\end{align*}
Maka $P = \{14, 35, 77, 46, 12\}$.

\textbf{2. Enkripsi Pesan} $M = (1, 0, 1, 0, 1)$:
\[ T = (1 \cdot 14) + (0 \cdot 35) + (1 \cdot 77) + (0 \cdot 46) + (1 \cdot 12) = 14 + 77 + 12 = 103 \]

\textbf{3. Dekripsi}:
\begin{itemize}
    \item Cari $m^{-1} \pmod{101}$: $7 \cdot 29 = 203 = 2 \cdot 101 + 1$. Jadi $m^{-1} = 29$.
    \item Hitung $T' = (103 \cdot 29) \pmod{101} = 2987 \pmod{101} = 58$.
    \item Selesaikan $T' = 58$ dengan $S = \{2, 5, 11, 21, 45\}$ menggunakan algoritma \textit{greedy}:
    \begin{itemize}
        \item $58 \ge 45 \rightarrow$ Ambil 45. Sisa $58-45 = 13$.
        \item $13 \ge 21$ (Tidak)
        \item $13 \ge 11 \rightarrow$ Ambil 11. Sisa $13-11 = 2$.
        \item $2 \ge 5$ (Tidak)
        \item $2 \ge 2 \rightarrow$ Ambil 2. Sisa $2-2 = 0$.
    \end{itemize}
    Hasil bit: $(1, 0, 1, 0, 1)$. (Terbukti!)
\end{itemize}

\subsection{Implementasi dalam Python}

\begin{pycode}
def modInverse(a, m):
    m0, y, x = m, 0, 1
    if m == 1: return 0
    while a > 1:
        q = a // m
        t = m
        m = a % m
        a = t
        y, x = x - q * y, y
    if x < 0: x = x + m0
    return x

def encrypt_knapsack(message, public_key):
    return sum(m_i * p_i for m_i, p_i in zip(message, public_key))

def decrypt_knapsack(ciphertext, private_key, m, q):
    S, m_inv = private_key, modInverse(m, q)
    t_prime = (ciphertext * m_inv) % q

    message = []
    for s in reversed(S):
        if t_prime >= s:
            message.append(1)
            t_prime -= s
        else:
            message.append(0)
    return message[::-1]

# Pengujian
S = [2, 5, 11, 21, 45]
m, q = 7, 101
P = [(s * m) % q for s in S]
msg = [1, 0, 1, 0, 1]
ct = encrypt_knapsack(msg, P)
pt = decrypt_knapsack(ct, S, m, q)
print(f"Msg: {msg} -> CT: {ct} -> PT: {pt}")
\end{pycode}

\begin{aktivitas}
  \item \textbf{Eksperimen Parameter}: Cobalah mengubah nilai $m$ dan $q$ pada implementasi Python di atas. Apa yang terjadi jika $\gcd(m, q) \neq 1$? Mengapa hal ini menyebabkan kegagalan dekripsi?
\end{aktivitas}

\begin{latihan}
  \item Diketahui kunci publik $P = \{15, 27, 51, 95, 187\}$ dan modulus $q=256$. Jika Anda mengetahui bahwa barisan superincreasing asal adalah $S = \{1, 2, 4, 8, 16\}$, tentukan nilai $m$ yang digunakan!
  \item Mengapa pemilihan $q$ harus lebih besar dari jumlah seluruh elemen barisan superincreasing $S$? Jelaskan konsekuensinya jika $q \le \sum s_i$!
  \item Buktikan bahwa jika $S$ adalah barisan superincreasing, maka setiap jumlah subset dari $S$ menghasilkan nilai yang unik!
\end{latihan}
